% Part: conditional-logics
% Chapter: minimal-change-semantics
% Section: true-false

\documentclass[../../../include/open-logic-section]{subfiles}

\begin{document}

\olfileid{con}{min}{tf}

\olsection{خلافِ واقعہ شرطیاں دی سچائی تے جھوٹ}

خلافِ واقعہ شرطیہ $!A \cif !B$ سچ اے (خالی شرط کارن نہیں) جے سبھ توں
نیڑلے $!A$ والے جہان سارے $!B$ والے جہان ہون، جیویں \olref{fig:true} وچ
وکھایا اے۔
\begin{figure}
\begin{center}
\begin{tikzpicture}[scale=.7]\small
  \spheresystem{5}
  \draw (0,0) node {$w$};
  \propositionintersect{0}{4}{40}{3.5}
  {\tikzset{proposition/.append={smooth,tension=1.4}}
  \proposition[shift={(1.6,-1)}]{90}{1}{30}{4}}
  \path (1.5,3) node[below] {$\formula{B}$};
  \path (3,0) node {$\formula{A}$};
\end{tikzpicture}
\caption{غیر خالی شرط نال سچا خلافِ واقعہ شرطیہ}
\ollabel{fig:true}
\end{center}
\end{figure}
خلافِ واقعہ شرطیہ $w$ اُتے اوہدوں وی سچ اے جے~$w$ دے آلے دوالے گولیاں
دے نظام وچ $!A$ نوں تھاں دین والا کوئی گولا ای نہیں۔ ایس صورت وچ اوہ
\emph{خالی شرط کارن} سچ اے (\olref{fig:vacuous} ویکھو)۔
\begin{figure}
\begin{center}
\begin{tikzpicture}[scale=.7]\small
  \spheresystem{5}
  \draw (0,0) node {$w$};
  \proposition{0}{6.5}{40}{3.5}
  {\tikzset{proposition/.append={smooth,tension=1.4}}
  \proposition[shift={(1.2,-1)}]{90}{1}{30}{4}}
  \path (1.5,3) node[below] {$\formula{B}$};
  \spherepos{0}{7}{node {$\formula{A}$}}
\end{tikzpicture}
\caption{خالی شرط کارن سچا خلافِ واقعہ شرطیہ}
\ollabel{fig:vacuous}
\end{center}
\end{figure}

ایہ دو طریقیاں نال جھوٹا ہو سکدا اے۔ اک صورت ایہ اے کہ سبھ توں نیڑلے
$!A$ والے جہاناں وچوں سارے $!B$ والے جہان نہ ہون، پر کجھ ہون۔ ایس صورت
وچ $!A \cif \lnot !B$ وی جھوٹا اے (\olref{fig:false} ویکھو)۔
\begin{figure}
\begin{center}
\begin{tikzpicture}[scale=.7]\small
  \spheresystem{5}
  \draw (0,0) node {$w$};
  \propositionintersect{0}{4}{40}{3.5}
  {\tikzset{proposition/.append={smooth,tension=1.4}}
  \proposition[shift={(1.6,0)}]{90}{1}{30}{3}}
  \path (1.5,3) node[below] {$\formula{B}$};
  \path (3,0) node {$\formula{A}$};
\end{tikzpicture}
\caption{جھوٹا خلافِ واقعہ شرطیہ، اُلٹی تالی والا وی جھوٹا}
\ollabel{fig:false}
\end{center}
\end{figure}
جے سبھ توں نیڑلے $!A$ والے جہاناں دا $!B$ والے جہاناں نال کوئی سانجھا
جہان نہیں، تاں $!A \cif !B$ جھوٹا اے۔ پر ایس صورت وچ سبھ توں نیڑلے
سارے $!A$ والے جہان $\lnot !B$ والے جہان نیں، سو $!A \cif \lnot !B$ سچ
اے (\olref{fig:false-opposite} ویکھو)۔
\begin{figure}
\begin{center}
\begin{tikzpicture}[scale=.7]\small
  \spheresystem{5}
  \draw (0,0) node {$w$};
  \propositionintersect{0}{4}{40}{3.5}
  {\tikzset{proposition/.append={smooth,tension=1.4}}
  \proposition[shift={(-1,0)}]{90}{1}{30}{3}}
  \path (-1,3) node[below] {$\formula{B}$};
  \path (3,0) node {$\formula{A}$};
\end{tikzpicture}
\caption{جھوٹا خلافِ واقعہ شرطیہ، اُلٹی تالی والا سچا}
\ollabel{fig:false-opposite}
\end{center}
\end{figure}

سخت شرطیے دے اُلٹ، خلافِ واقعہ شرطیے اتفاقی طور اُتے سچ یا جھوٹ ہو
سکدے نیں۔ \olref{fig:contingent} والا گولیاں دا نمونہ ویکھو۔
$!A$ والے جہان جیہڑے~$u$ دے سبھ توں نیڑے نیں، اوہ سارے $!B$ والے جہان نیں، سو
$\mSat{M}{!A \cif !B}[u]$۔ پر $!A$ والے جہان جیہڑے~$v$ دے سبھ توں نیڑے نیں، اوہناں وچوں کجھ جہان
$!B$ والے جہان نہیں، سو $\mSat/{M}{!A \cif !B}[v]$۔
\begin{figure}
\begin{center}
\begin{tikzpicture}[scale=.6]\small
  \spheresystem{7}
  \spheresystem[shift={(0:3.1)},dashed]{4}
  \propositionintersect{45}{4}{40}{4.5}
  \begin{scope}
    \clip \propositionplot{45}{4}{40}{4.5} ;
    \spherelayer[shift={(0:3.1)}]{3}
  \end{scope}
  \proposition[shift={(1.4,-.5)}]{90}{1}{35}{5}
  \path (0,0) node {$u$};
  \path (0:3.1) node {$v$};
  \spherepos{35}{9}{node {$\formula{A}$}}
  \spherepos{80}{9}{node {$\formula{B}$}}
\end{tikzpicture}
\end{center}
\caption{اتفاقی خلافِ واقعہ شرطیہ}
\ollabel{fig:contingent}
\end{figure}
\end{document}
